%% test-basic.tex
%% ==============
%% Typesetting test for every engine: the footnotes of the boxes with
%% "page footnote" are page footnotes on the page of their marks.
%%
%%   pdflatex test-basic.tex   (or xelatex / lualatex)
%%   platex  test-basic.tex && dvipdfmx test-basic
%%   uplatex -jobname=test-basic "\def\CLS{ujarticle}\input test-basic"
%% The class is article by default; set \CLS for another one (jarticle,
%% ujarticle, ltjsarticle, ...).  With pTeX the class option dvipdfmx
%% is given.
%%
%% What is covered:
%%   * footnotes in the body text, and in the title, upper and lower
%%     part of a box, several in one box;
%%   * \footnote[<num>], \footnotemark + \footnotetext in a box;
%%   * a box with "page footnote" inside another one, and an ordinary
%%     box (a footnote of the box) inside a box with "page footnote";
%%   * an ordinary box for comparison;
%%   * \tcbox, a box defined by \newtcolorbox, and \tcbset{page footnote};
%%   * a box that does not fit at the bottom of a page and moves to the
%%     next page: its footnotes move with it.
%%
%% test-hyperref.tex runs this file with hyperref.
%%
%% What to check:
%%   * TPF-RESULT OK, TPF-COUNTER OK, and no error or warning.
\ifdefined\CLS\else\def\CLS{article}\fi
\ifdefined\ptexversion\def\CLSOPT{dvipdfmx}\else\def\CLSOPT{}\fi
\expandafter\documentclass\expandafter[\CLSOPT]{\CLS}
\usepackage{tcbpagefootnote}
\ifdefined\TPFhyperref\usepackage{hyperref}\fi
%% test-common.tex
%% ===============
%% Helpers shared by the test files; \input this file after
%% \usepackage{tcbpagefootnote}.
%%
%% \TPFfootnote[<kind>]{<text>}
%%   \footnote{<text>} with a check of where its text is typeset.
%%   <kind> (default: page) is one of
%%     page   a page footnote on the page of its mark;
%%     box    a footnote of the box (as in a minipage) on the page of
%%            its mark;
%%     later  a page footnote on a later page than its mark (a box
%%            broken across pages).
%% \TPFmark{<id>} and \TPFtext{<id>}{<kind>}
%%   The two halves of \TPFfootnote, for \footnote[<num>],
%%   \footnotemark and \footnotetext: put \TPFmark right after the mark
%%   and \TPFtext at the beginning of the text.  <id> is any name that
%%   is not a number (numbers are used by \TPFfootnote).
%% \TPFcounter{<counter>}{<value>}
%%   Check the value of <counter> at this point.
%%
%% For each mark and each text, the page, the mark (\@thefnmark) and,
%% for a text, whether it is set in the full \columnwidth of the page
%% (a page footnote) are written to the .aux file at shipout.  When the
%% .aux file is read back at \end{document}, every mark is checked:
%%   * its text is typeset exactly once, with the same mark;
%%   * the text is where <kind> says.
%% Logs "TPF-NOTE NG: ..." for each failure, then "TPF-RESULT OK (<n>
%% notes)" or "TPF-RESULT NG".  A single run is enough.
\makeatletter
\ExplSyntaxOn
\int_new:N  \g__tpf_id_int
\seq_new:N  \g__tpf_ids_seq
\bool_new:N \g__tpf_ng_bool
\dim_new:N  \g__tpf_colwd_dim
\tl_new:N   \l__tpf_tmp_tl
\AddToHook { begindocument/end }
  { \dim_gset:Nn \g__tpf_colwd_dim { \columnwidth } }

\cs_new_protected:Npn \__tpf_ng:n #1
  {
    \bool_gset_true:N \g__tpf_ng_bool
    \iow_term:e { TPF-NOTE~NG:~#1 }
  }

%% Write to the .aux file at shipout; #1 is frozen now except \thepage
\cs_new_protected:Npn \__tpf_shipout:n #1
  {
    \tl_set:Ne \l__tpf_tmp_tl {#1}
    \exp_args:NNV \iow_shipout_e:Nn \@auxout \l__tpf_tmp_tl
  }
\cs_new:Npn \__tpf_frozen_mark:
  { \exp_not:N \exp_not:n { \exp_not:o { \@thefnmark } } }

\NewDocumentCommand \TPFmark { m }
  {
    \seq_gput_right:Nn \g__tpf_ids_seq {#1}
    \__tpf_shipout:n
      {
        \exp_not:N \exp_not:N \exp_not:N \TPF@m
        {#1} { \exp_not:N \thepage } { \__tpf_frozen_mark: }
      }
  }
\NewDocumentCommand \TPFtext { m m }
  {
    \__tpf_shipout:n
      {
        \exp_not:N \exp_not:N \exp_not:N \TPF@t
        {#1} {#2} { \exp_not:N \thepage } { \__tpf_frozen_mark: }
        { \dim_compare:nNnTF { \hsize } = { \g__tpf_colwd_dim } { 1 } { 0 } }
      }
  }
\NewDocumentCommand \TPFfootnote { O{page} +m }
  {
    \int_gincr:N \g__tpf_id_int
    \exp_args:NV \__tpf_footnote:nnn \g__tpf_id_int {#1} {#2}
  }
\cs_new_protected:Npn \__tpf_footnote:nnn #1#2#3
  { \footnote { \TPFtext {#1} {#2} #3 } \TPFmark {#1} }

\NewDocumentCommand \TPFcounter { m m }
  {
    \int_compare:nNnTF { \value {#1} } = {#2}
      { \iow_term:e { TPF-COUNTER~OK:~#1=#2 } }
      { \__tpf_ng:n { counter~#1=\int_use:c { c@#1 },~want~#2 } }
  }

%% The records in the .aux file are ignored when it is read at
%% \begin{document} (they come from the previous run)
\cs_new_eq:NN \TPF@m \use_none:nnn
\cs_new_eq:NN \TPF@t \use_none:nnnnn
\cs_new_protected:Npn \__tpf_record_mark:nnn #1#2#3
  {
    \tl_gset:cn { g__tpf_mpage_#1_tl } {#2}
    \tl_gset:cn { g__tpf_mmark_#1_tl } {#3}
  }
\cs_new_protected:Npn \__tpf_record_text:nnnnn #1#2#3#4#5
  {
    \int_if_exist:cF { g__tpf_tcount_#1_int }
      { \int_new:c { g__tpf_tcount_#1_int } }
    \int_gincr:c { g__tpf_tcount_#1_int }
    \tl_gset:cn { g__tpf_tkind_#1_tl } {#2}
    \tl_gset:cn { g__tpf_tpage_#1_tl } {#3}
    \tl_gset:cn { g__tpf_tmark_#1_tl } {#4}
    \tl_gset:cn { g__tpf_twide_#1_tl } {#5}
  }
\AddToHook { enddocument/afterlastpage }
  {
    \cs_gset_eq:NN \TPF@m \__tpf_record_mark:nnn
    \cs_gset_eq:NN \TPF@t \__tpf_record_text:nnnnn
  }
\AddToHook { enddocument/afteraux } { \__tpf_check: }

\cs_new_protected:Npn \__tpf_check:
  {
    \seq_map_function:NN \g__tpf_ids_seq \__tpf_check_one:n
    \seq_if_empty:NT \g__tpf_ids_seq { \__tpf_ng:n { no~footnotes } }
    \bool_if:NTF \g__tpf_ng_bool
      { \iow_term:n { TPF-RESULT~NG } }
      {
        \iow_term:e
          { TPF-RESULT~OK~(\seq_count:N \g__tpf_ids_seq\c_space_tl notes) }
      }
  }
\cs_new_protected:Npn \__tpf_check_one:n #1
  {
    \tl_if_exist:cTF { g__tpf_mpage_#1_tl }
      {
        \int_if_exist:cTF { g__tpf_tcount_#1_int }
          {
            \int_compare:nNnTF { \int_use:c { g__tpf_tcount_#1_int } } = { 1 }
              { \__tpf_check_text:n {#1} }
              {
                \__tpf_ng:n
                  { #1:~text~typeset~\int_use:c { g__tpf_tcount_#1_int }~times }
              }
          }
          { \__tpf_ng:n { #1:~text~lost } }
      }
      { \__tpf_ng:n { #1:~mark~not~shipped~out } }
  }
\cs_new_protected:Npn \__tpf_check_text:n #1
  {
    \tl_set:Ne \l__tpf_tmp_tl
      {
        #1:~\tl_use:c { g__tpf_tkind_#1_tl },~
        mark~\tl_use:c { g__tpf_mmark_#1_tl }~
        on~p.\tl_use:c { g__tpf_mpage_#1_tl },~
        text~\tl_use:c { g__tpf_tmark_#1_tl }~
        on~p.\tl_use:c { g__tpf_tpage_#1_tl }~
        \str_if_eq:eeTF { \tl_use:c { g__tpf_twide_#1_tl } } { 1 }
          { (page) } { (box) }
      }
    \bool_if:nTF
      {
        \tl_if_eq_p:cc { g__tpf_mmark_#1_tl } { g__tpf_tmark_#1_tl }
        && \__tpf_where_p:n {#1}
      }
      { \iow_term:e { TPF-NOTE~OK:~\l__tpf_tmp_tl } }
      { \__tpf_ng:n { \l__tpf_tmp_tl } }
  }
\prg_new_conditional:Npnn \__tpf_where:n #1 { p }
  {
    \str_case_e:nnF { \tl_use:c { g__tpf_tkind_#1_tl } }
      {
        { page }
          {
            \bool_if:nTF
              {
                \int_compare_p:nNn { \tl_use:c { g__tpf_tpage_#1_tl } }
                  = { \tl_use:c { g__tpf_mpage_#1_tl } }
                && \str_if_eq_p:ee { \tl_use:c { g__tpf_twide_#1_tl } } { 1 }
              }
              { \prg_return_true: } { \prg_return_false: }
          }
        { box }
          {
            \bool_if:nTF
              {
                \int_compare_p:nNn { \tl_use:c { g__tpf_tpage_#1_tl } }
                  = { \tl_use:c { g__tpf_mpage_#1_tl } }
                && \str_if_eq_p:ee { \tl_use:c { g__tpf_twide_#1_tl } } { 0 }
              }
              { \prg_return_true: } { \prg_return_false: }
          }
        { later }
          {
            \bool_if:nTF
              {
                \int_compare_p:nNn { \tl_use:c { g__tpf_tpage_#1_tl } }
                  > { \tl_use:c { g__tpf_mpage_#1_tl } }
                && \str_if_eq_p:ee { \tl_use:c { g__tpf_twide_#1_tl } } { 1 }
              }
              { \prg_return_true: } { \prg_return_false: }
          }
      }
      { \prg_return_false: }
  }
\ExplSyntaxOff

%% \TPFfiller{<n>}: a paragraph of <n> filler sentences
\newcommand\TPF@filler{%
  This is a filler sentence of plain English text for the test. }
\newcommand\TPFfiller[1]{%
  \count@=0
  \loop\advance\count@ by 1 \TPF@filler\ifnum\count@<#1 \repeat\par}
\makeatother
\endinput

\newtcolorbox{notebox}[1][]{colback=blue!5!white, page footnote, #1}
\begin{document}
\section*{Footnotes in boxes}

Body text\TPFfootnote{A footnote in the body text.}.

\begin{tcolorbox}[title=Title\TPFfootnote{A footnote in the title.},
  page footnote]
Upper part\TPFfootnote{The first footnote in the upper part.} and
again\TPFfootnote{The second footnote in the upper part.}.
\tcblower
Lower part\TPFfootnote{A footnote in the lower part.}.
\end{tcolorbox}

An ordinary box for comparison:
\begin{tcolorbox}
Ordinary box\TPFfootnote[box]{A footnote of the ordinary box.}.
\end{tcolorbox}

\begin{tcolorbox}[page footnote]
Optional argument\footnote[9]{\TPFtext{opt}{page}A footnote numbered
9.}\TPFmark{opt}, and footnotemark\footnotemark\TPFmark{fm}.
\footnotetext{\TPFtext{fm}{page}The text for the footnotemark.}
\end{tcolorbox}

Body text after the box\TPFfootnote{The numbering goes on.}.

\newpage
\section*{Nested boxes}

\begin{tcolorbox}[page footnote]
Outer box\TPFfootnote{The first footnote of the outer box.}.
\begin{tcolorbox}[page footnote]
Inner box\TPFfootnote{A footnote of the inner box.}.
\end{tcolorbox}
Outer box again\TPFfootnote{The second footnote of the outer box.}.
\begin{tcolorbox}
Ordinary inner box\TPFfootnote[box]{A footnote of the ordinary inner
  box.}.
\end{tcolorbox}
\end{tcolorbox}

\section*{Other boxes}

Inline \tcbox[page footnote]{tcbox\TPFfootnote{A footnote in a tcbox.}}
in a paragraph.

\begin{notebox}[title=notebox]
Box by newtcolorbox\TPFfootnote{A footnote in a notebox.}.
\end{notebox}

\begingroup
\tcbset{page footnote}
\begin{tcolorbox}
Set by tcbset\TPFfootnote{A footnote in a box set by tcbset.}.
\end{tcolorbox}
\endgroup

\newpage
\section*{A box at the bottom of a page}

\TPFfiller{1}
\vspace{\dimexpr\textheight-12\baselineskip\relax}
\TPFfiller{1}
\begin{tcolorbox}[page footnote]
This box does not fit at the bottom of the page\TPFfootnote{A footnote of
  the box moved to the next page.}.
\TPFfiller{6}
\end{tcolorbox}

Body text after the box\TPFfootnote{The last footnote.}.

\TPFcounter{footnote}{15}
\TPFcounter{mpfootnote}{0}
\end{document}
